\chapter{Example: electron spin resonance of the doublet}\label{ch:ex-esr}

In a static magnetic field the doublet splits into $D_\up$ and $D_\dn$. A transverse ac field
then drives spin flips between them: electron spin resonance (ESR) of the impurity spin, the
mechanism behind spin qubits in quantum dots~\cite{koppens2006driven,nowack2007coherent,golovach2006electricdipoleinduced}
and ESR-STM of single atoms~\cite{baumann2015electron,willke2018probing,yang2019coherent}. In
a superconducting junction the spin is partially screened by the coupling to the leads, and
ESR probes the resulting reduction of the local spin response. This renormalization
is also relevant to Andreev spin qubits, whose electrically driven transitions
depend on spin--orbit coupling~\cite{chtchelkatchev2003andreev,park2017andreev,hays2021coherent,pitavidal2023direct,bargerbos2023spectroscopy}.

\section{The model}

The impurity carries a static Zeeman energy $B\Sz_d$ and a transverse drive,
\begin{equation}
  \Hd(t)=\frac U2(n_d-1)^2+\eps(n_d-1)+B\,\Sz_d+b\cos(\Omega t)\,S^x_d,\qquad
  S^x_d=\tfrac12\bigl(d_\up\dg d_\dn+d_\dn\dg d_\up\bigr).
\end{equation}
The field acts on the dot only, as for a $g$-factor that is large on the dot and small in the
leads. The transverse term couples different $S_z$ values, so the calculation is done in the odd
sector with all spin projections, \code{Sector(1, None)}. A model constructor with a
\code{transverse} parameter implements this drive:

\begin{lstlisting}
from qdjj_solver import Impurity, annihilate, create, hermitian_pair, make_model

def esr_model(transverse=0., gamma=.3, u=1., detuning=0., field=.02, phi=.8):
    anderson = Impurity.anderson(u, detuning, field)
    operator = anderson.operator + hermitian_pair(.5*transverse*create(0)*annihilate(1))  # b S_x
    impurity = Impurity(1, operator, metadata=dict(kind="anderson-transverse"))
    t = np.sqrt(gamma/(np.pi*bath.rho*2))*np.eye(2)
    return make_model(impurity, [bath, bath], [t, t], phases=[-phi/2, phi/2],
                      phase_velocities=[-.5, .5])

esr = PeriodicHamiltonian.from_parameters(esr_model, dict(transverse=0., gamma=.3),
                                          ParameterDrive.cosine("transverse", .004), omega=.0173)
\end{lstlisting}

\section{Rabi splitting and the effective \texorpdfstring{$g$}{g}-factor}

For a free spin, the Zeeman splitting is $B$ and the resonant Rabi splitting of the drive
$b\cos(\Omega t)S^x$ is $b/2$. In the junction the doublet is a correlated state of the dot
spin and the leads. Within the doublet the impurity spin operator is proportional to the
total spin, $\bm S_d=(g_{\mathrm{eff}}/g)\,\bm S$, by the Wigner--Eckart theorem, with
\begin{equation}\label{eq:geff}
  \frac{g_{\mathrm{eff}}}{g}=2\abs{\matel{D_\dn}{\Sz_d}{D_\dn}}\le1 .
\end{equation}
The relation is exact at zero field. In a field the polarizations of the two Zeeman levels
differ at linear order in $B$; we quote that of the lower level $D_\dn$, and the transverse
element $2\abs{\matel{D_\up}{S^x_d}{D_\dn}}$ equals their mean.
Both the Zeeman splitting and the Rabi splitting are therefore renormalized by the same
factor:
\begin{equation}\label{eq:esr-renorm}
  E_{D_\up}-E_{D_\dn}\simeq\frac{g_{\mathrm{eff}}}g\,B,\qquad
  \delta_{\mathrm{Rabi}}=b\,\abs{\matel{D_\up}{S^x_d}{D_\dn}}=\frac{g_{\mathrm{eff}}}g\,\frac b2 .
\end{equation}
The reduction of the effective $g$-factor of a quantum-dot junction with increasing coupling
is the impurity Knight shift~\cite{pavesic2023impurity}.

\Cref{fig:esr}(a) shows the Autler--Townes splitting of the doublet at the ESR resonance
for $\Gamma=\q{esr.gamma}$, where $g_{\mathrm{eff}}/g=\q{esr.geff}$. Panel (b) compares three
determinations of $g_{\mathrm{eff}}/g$ as functions of $\Gamma$: the static spin polarization
\eqref{eq:geff}, the Zeeman splitting divided by $B$, and the Floquet Rabi splitting divided
by $b/2$. They agree within $\q{esr.agreement}$, the size of the finite-field and Bloch--Siegert
corrections, and decrease from $\q{esr.geff_first}$ at $\Gamma=\q{esr.gamma_first}$ to
$\q{esr.geff_last}$ at $\Gamma=\q{esr.gamma_last}$.

\begin{figure}[htb]
\centering
\manualfigure{esr}
\caption{ESR of the doublet in the reference junction at $\eps=\q{esr.epsilon}$, $U=\q{esr.u}$,
$\phi=\q{esr.phi}$, $B=\q{esr.field}$. (a) Rotating-frame quasienergies near the Zeeman
resonance for $\Gamma=\q{esr.gamma}$ and $b=\q{esr.amplitude}$, Sambe (solid) and RWA
(dashed). (b) Effective $g$-factor from the static spin polarization, the Zeeman splitting
and the Floquet Rabi splitting.}
\label{fig:esr}
\end{figure}

\section{Probing the driven spin}\label{sec:esr-probe}

The dressed doublet can be probed with a second, weak field (\cref{ch:probes}). A transverse
probe $S^x_d$ acts on the pumped transition itself, and a longitudinal probe $\Sz_d$, a weak
modulation of the static field, connects the two dressed states. \Cref{fig:esr-probe} shows
both for $\Gamma=\q{esr_probe.gamma}$, $b=\q{esr_probe.amplitude}$, starting from the dressed
continuation of the lower Zeeman level. The Zeeman splitting is $Z=\q{esr_probe.zeeman}$ and
the resonant Rabi splitting is $\delta=\q{esr_probe.rabi}=\q{esr_probe.ratio}\,Z$.

The transverse probe, panels (a) and (b), shows the Mollow pattern of \cref{sec:mollow}:
two lines that avoid each other at resonance and follow the rotating-wave prediction
$\Omega\pm w$, with $w=\sqrt{(\Omega-Z)^2+\delta^2}$, and the elastic line at $\Omega$. By the
Wigner--Eckart argument above, each of the four lines carries
$\abs{\matel{D_\up}{S^x_d}{D_\dn}}^2/4\simeq(g_{\mathrm{eff}}/g)^2/16=\q{esr_probe.sx_rwa}$ at
resonance; the computed values are $\q{esr_probe.sx_lower}$ and $\q{esr_probe.sx_upper}$ for
the sidebands and $\q{esr_probe.sx_elastic}$ for each elastic line. The asymmetry of the two
sidebands is a counter-rotating effect of order $\delta/Z$: it halves when $b$ is halved.

The longitudinal probe, panel (c), shows a single line at the Rabi frequency, the transition between the two dressed
states. Its strength,
$\tfrac14\bigl(\matel{D_\up}{\Sz_d}{D_\up}-\matel{D_\dn}{\Sz_d}{D_\dn}\bigr)^2\simeq(g_{\mathrm{eff}}/g)^2/4
=\q{esr_probe.sz_rwa}$ in the RWA (computed: $\q{esr_probe.sz_line}$), measures the effective
$g$-factor through a matrix element instead of a splitting.

\begin{figure}[htb]
\centering
\manualfigure{esr_probe}
\caption{Probing the Rabi-driven doublet ($\Gamma=\q{esr_probe.gamma}$,
$b=\q{esr_probe.amplitude}$, $M=\q{esr_probe.harmonics}$, static basis
$K=\q{esr_probe.basis}$). (a) Lines of a transverse probe $S^x_d$ versus drive frequency,
darker for stronger lines; dashed: RWA $\Omega\pm w$. (b) Rate spectrum of $S^x_d$ at
resonance: sidebands at $Z\pm\delta$ and the central elastic line. (c) Rate spectrum of
$\Sz_d$ at resonance: the line at the Rabi frequency $\delta=\q{esr_probe.ratio}\,Z$. Line
half width $\q{esr_probe.width}\,Z$.}
\label{fig:esr-probe}
\end{figure}

\begin{caution}
At resonance the replicas $\ket{D_\dn,m}$ and $\ket{D_\up,m-1}$ form a resonant chain through
the entire Sambe space. A small Floquet cutoff truncates this chain and creates unpaired edge
states, which perturb the interior replicas at the order of the Rabi coupling. Even for weak
drives, use a Floquet cutoff of several harmonics ($M\gtrsim5$); the residual bound and the
check that the tracked states are distinct Floquet states detect an insufficient cutoff.
\end{caution}

\tutorial{10_spin_resonance}
